\section{Conclusion}\label{sec:conclusion}

A no-trade clause is a promise made by a team to a player. A long-term contract is a promise made by a player to a team, but only the first promise is enforceable. When the second fails, the first becomes a weapon. The protected player can name his buyers, signal that the seller must sell, and, by going public, destroy the seller's last threat. The surplus that changes hands goes not to the player but to the team he chooses.

The data show that the NHL has moved steadily into the region where this matters. Nearly half of all cap dollars, more than three-quarters of star contracts, and a rising share of future commitments now carry a veto. Protection keeps players in place, roughly halving the annual chance that a protected player is traded. But a growing share of trades now pass through those vetoes, and protection has not been priced. The theory predicts the response: narrower vetoes, decaying clauses, explicit prices, and a fight at the next bargaining table. Whether the 2026 wave produces that response is the natural test of the argument, and the data assembled here make it possible to run.
